\section{引言}\label{sec:introduction}

大数定律（Law of Large Numbers）是概率论中历史最悠久、地位最根本的极限定理。它的经验原型人所共知：\textbf{大量重复试验中，事件发生的频率会稳定在其概率附近}；掷一枚均匀硬币，掷出的次数越多，正面出现的比例越接近 $1/2$。把这一直觉精确化，便得到定律的经典形式：

\begin{quote}
    \textit{设 $X_1, X_2, \dots$ 为独立同分布的随机变量序列，公共分布的数学期望 $\mu$ 有限，$S_n = X_1 + \cdots + X_n$ 为部分和，则经验平均 $S_n/n$ 依概率（弱大数定律）或几乎必然（强大数定律）收敛于 $\mu$。}
\end{quote}

这一定律的深刻之处在于：它把\textbf{不可预测的单次试验}与\textbf{可预测的集体行为}连接起来——随机性在个体层面不可约减，却在集体层面服从精确的数学规律。正如伯努利在《猜度术》中所说，即使最愚钝的人凭本能也知晓这一规律，但\textbf{以数学证明它}却是一项艰巨的事业：伯努利为此苦心经营二十余年，并感叹其难度远超几何学中的任何命题。

\subsection{弱形式与强形式}

大数定律有强弱两种形式，其间的差别直到二十世纪初才被真正理解：

\begin{itemize}
    \item \textbf{弱大数定律}（WLLN）：$S_n/n \xrightarrow{P} \mu$，即对任意 $\varepsilon > 0$，$\lim_{n\to\infty} P\left(\left|S_n/n - \mu\right| > \varepsilon\right) = 0$。这是对\textbf{每个固定时刻} $n$ 的收敛性陈述。
    \item \textbf{强大数定律}（SLLN）：$S_n/n \to \mu$ 几乎必然（a.s.），即整条样本路径上的收敛同时成立。这是对\textbf{整条轨道}的陈述，蕴含弱形式但远强于之。
\end{itemize}

两种形式的逻辑关系并不平凡：几乎必然收敛蕴含依概率收敛，反之不然；但若把序列换成子列，依概率收敛总能抽出几乎必然收敛的子列。历史上，弱形式（伯努利、切比雪夫）先于强形式（博雷尔、康泰利、科尔莫哥洛夫）出现约两百年，原因在于强形式需要对无穷多个时刻的一致控制——这正是科尔莫哥洛夫极大值不等式所要克服的核心困难。

\subsection{研究意义}

大数定律的研究意义体现在多重层面：

\begin{enumerate}[label=(\arabic*)]
    \item \textbf{概率论的奠基}：大数定律是\textbf{第一极限定理}，它使“概率”从主观信念或频率的循环定义中解放出来——频率趋近概率本身成为一个可证明的数学定理，为概率论成为严格的数学学科（1933 年科尔莫哥洛夫公理化\cite{kolmogorov1933}）提供了第一块基石。
    \item \textbf{统计推断的逻辑基础}：统计学的根本信条“样本信息可以推断总体”之所以成立，正因经验分布函数一致收敛于总体分布（Glivenko--Cantelli 定理\cite{glivenko1933}）、相合估计量的存在性等，而这些都是大数定律的直接推论。
    \item \textbf{计算随机化的理论保障}：蒙特卡洛方法以随机模拟代替确定性计算，其误差以 $O(n^{-1/2})$ 的速率下降——这一保障来自大数定律（与中心极限定理的联合运用）。
    \item \textbf{推广范式的原型}：从独立到两两独立（Etemadi\cite{etemadi1981}）、到鞅差（杜布）、到平稳遍历（伯克霍夫遍历定理\cite{birkhoff1931}），大数定律的每一次条件弱化都是概率论技术的试金石，其推广史本身就是一部现代概率论方法史。
\end{enumerate}

\subsection{本文结构}

本文组织如下：\cref{sec:history} 回顾从《猜度术》到现代的历史发展；\cref{sec:methods} 介绍主要证明方法；\cref{sec:results} 总结重要研究成果；\cref{sec:applications} 讨论相关应用与联系；\cref{sec:open_problems} 列举开放问题与推广方向；\cref{sec:conclusion} 总结全文。
